\section{Experiments}
\label{sec:experiments}

\subsection{Experimental Setup}

The navigation system is evaluated in two distinct environments to assess its performance under different conditions: a simulated environment using Gazebo and a real-world deployment on a TurtleBot3 Burger platform.

\subsubsection{Simulation Environment}
The simulation experiments are conducted using:
\begin{itemize}
    \item Gazebo with TurtleBot3 Burger model
    \item Simulated LiDAR sensor (360° field of view)
    \item Static obstacle configurations
    \item Ground-truth odometry for pose estimation
    \item Dynamic goal generation through odometry messages
\end{itemize}


\subsubsection{Real Robot Platform}
Physical experiments are performed using:
\begin{itemize}
    \item TurtleBot3 Burger robot 
    \item AprilTag as the moving target
    \item Static and dynamic obstacle configurations
\end{itemize}

The real robot experiments are conducted in a laboratory environment with movable obstacles, the target is moved manually while holding an AprilTag marker, creating dynamic goal positions for the robot to follow.

\subsection{Evaluation Metrics}
\label{ssec:metrics}
To quantitatively assess the navigation system performance, the following metrics are recorded during experimental trials:

\begin{itemize}
    \item \textbf{Success rate}: Percentage of trials where the robot successfully reached the goal without collisions or timeouts
    \item \textbf{Time of tracking}: amount of time during which the robot did not lose the target. 
    \item \textbf{Root Mean Square Error (RMSE)} of the target distance and bearing from the target moving robot, considering as optimal value the distance we define in the objective function of the DWA and 0 for the bearing angle. 
    \item \textbf{Average and Minimum distance to obstacles} using the /scan LiDAR data. 
\end{itemize}

\subsection{Parameter Configuration}

The DWA controller is configured with the parameters shown in \autoref{tab:dwa_params}, tuned through iterative testing in simulation and adapted for the real robot platform.

\begin{table}[h]
\centering
\caption{DWA Controller Parameters}
\label{tab:dwa_params}
\begin{tabular}{lccc}
\hline
\textbf{Parameter} & \textbf{Simulation} & \textbf{Real} & \textbf{Unit} \\
\hline
Control frequency & 15 & 15 & Hz \\
Simulation time& 2.0 & 2.0 & s \\
Time granularity & 0.1 & 0.1 & s \\
Linear velocity samples & 20 & 20 & - \\
Angular velocity samples & 20 & 20 & - \\
LiDAR angular sectors ($N_{ranges}$) & 20 & 20 & - \\
Goal distance tolerance & 0.2 & 0.35 & m \\
Collision tolerance (obj) & 0.5 & 0.5 & m \\
Collision tolerance (stop) & 0.15 & 0.18 & m \\
Max linear velocity & 0.22 & 0.22 & m/s \\
Max angular velocity & 2.84 & 2.84 & rad/s \\
Max linear acceleration & 0.22 & 0.22 & m/s² \\
Max angular acceleration & $\pi$ & $\pi$ & rad/s² \\
Robot radius & 0.2 & 0.2 & m \\
\hline
\multicolumn{4}{c}{\textit{Objective function weights}} \\
\hline
$\alpha$ (heading) & 0.1 & 0.1 & - \\
$\beta$ (velocity) & 0.2 & 0.2 & - \\
$\gamma$ (obstacle) & 0.2 & 0.2 & - \\
\hline
\multicolumn{4}{c}{\textit{Task-specific parameters}} \\
\hline
Slowdown threshold (Task 2a) & 0.25 & 0.25 & m \\
Slowdown weight (Task 2a) & 0.2 & 0.1 & - \\
Target distance (Task 2b) & 0.2 & 0.3 & m \\
$\delta$ (target distance weight, Task 2b) & 0.3 & 0.1 & - \\
\hline
\end{tabular}
\end{table}

\begin{table}[h]
\centering
\caption{Metrics}
\label{tab:metrics}
\begin{tabular}{lcccccccc}
\hline
\textbf{Metric} & \textbf{s1} & \textbf{s2a} & \textbf{s2b} & \textbf{r1} & \textbf{r1\_static} & \textbf{r1\_timeout} & \textbf{r2a} & \textbf{r2b}\\
\hline
Success Rate & 100.000\% & 100.000\% & 100.000\% & 42.353\% & 71.196\% & 0.000\% & 64.352\% & 4.082\%\\
Collision Rate & 0.000\% & 0.000\% & 0.000\% & 57.647\% & 0.000\% & 0.000\% & 35.648\% & 95.918\%\\
Timeout Rate & 0.000\% & 0.000\% & 0.000\% & 0.000\% & 28.804\% & 100.000\% & 0.000\% & 0.000\%\\
Tracking Percentage & 100.000 & 100.000 & 100.000 & 73.176 & 59.471 & 46.065 & 66.535 & 76.913\\
Rmse Distance & 0.313 & 0.316 & 0.320 & 1.669 & 2.685 & 1.222 & 1.223 & 1.253\\
Rmse Bearing & 0.133 & 0.132 & 0.145 & 0.224 & 0.107 & 0.071 & 0.186 & 0.174\\
Avg Distance & 0.431 & 0.444 & 0.448 & 1.634 & 2.550 & 1.407 & 1.307 & 1.418\\
Avg Bearing & 0.074 & 0.079 & 0.086 & 0.178 & 0.071 & 0.059 & 0.146 & 0.138\\
Avg Dist Obst & 1.9045847654342651 & 1.8823645114898682 & 1.8795063495635986 & 2.075275421142578 & 2.2188193798065186 & 1.9734026193618774 & 2.0853466987609863 & 2.4157567024230957\\
Min Dist Obst & 0.17209821939468384 & 0.1819748878479004 & 0.17993928492069244 & 0.1550000011920929 & 0.1979999989271164 & 0.30399999022483826 & 0.1550000011920929 & 0.1550000011920929\\
\hline
\end{tabular}
\end{table}

\paragraph{Note on tolerance parameters:} The goal distance tolerance is used both in the objective function evaluation and as the actual termination condition to verify goal achievement. For obstacle avoidance, two distinct tolerances are employed: a higher value (collision tolerance for obj) is used in the objective function to discourage proximity to obstacles, while a lower value (collision tolerance for stop) serves as the safety threshold that triggers an immediate stop in case of imminent collision.
